% Documento autónomo, figuras vectoriales y bibliografía integradas.
% Generado desde resultados verificados del experimento v2.0.0.
% Compilar con pdfLaTeX; repetir para actualizar índice y referencias.
% TFM completo: capítulos 1 a 7 y anexos A a D. Experimento reproducible adjunto.
\documentclass[11pt,a4paper,oneside]{report}
\usepackage[utf8]{inputenc}
\usepackage[T1]{fontenc}
\usepackage[spanish,es-tabla]{babel}
\usepackage[top=2.5cm,bottom=2.5cm,left=3cm,right=3cm]{geometry}
\usepackage{lmodern}
\usepackage{microtype}
\usepackage{setspace}
\usepackage{amsmath,amssymb}
\usepackage{booktabs,tabularx,array,longtable}
\usepackage{etoolbox}
\usepackage{pgfplots}
\usetikzlibrary{positioning,arrows.meta}
\pgfplotsset{compat=1.18}
\definecolor{modeloLR}{RGB}{36,110,152}
\definecolor{modeloSVM}{RGB}{184,103,34}
\usepackage{enumitem}
\usepackage{titlesec}
\usepackage{fancyhdr}
\usepackage{xurl}
\usepackage[hidelinks]{hyperref}
\hypersetup{pdftitle={Evaluación y mitigación de sesgos en modelos de inteligencia artificial aplicados al reclutamiento mediante currículums sintéticos},pdfauthor={Jordan Rajo Vázquez},pdflang={es-ES}}
\setstretch{1.04}
\setlength{\parindent}{1.2em}
\setlength{\parskip}{0.22em}
\setlength{\emergencystretch}{2em}
\setcounter{tocdepth}{1}
\setcounter{secnumdepth}{2}
\setlist{itemsep=0.25em,topsep=0.4em}
\titleformat{\chapter}[hang]{\normalfont\LARGE\bfseries}{\thechapter.}{0.6em}{}
\titlespacing*{\chapter}{0pt}{0pt}{18pt}
\titleformat{\section}{\normalfont\large\bfseries}{\thesection}{0.6em}{}
\titleformat{\subsection}{\normalfont\normalsize\bfseries}{\thesubsection}{0.6em}{}
\pagestyle{fancy}
\fancyhf{}
\fancyfoot[C]{\thepage}
\renewcommand{\headrulewidth}{0pt}
\fancypagestyle{plain}{\fancyhf{}\fancyfoot[C]{\thepage}\renewcommand{\headrulewidth}{0pt}}
\addto\captionsspanish{\renewcommand{\bibname}{Bibliografía}}
% Numeración del capítulo de bibliografía sin duplicar la entrada en el índice.
\patchcmd{\thebibliography}{\chapter*{\bibname}}{\chapter{\bibname}}{}{\PackageError{TFM}{No se pudo numerar la bibliografia}{Revisar la definicion de thebibliography.}}

\newcommand{\tituloTFM}{Evaluación y mitigación de sesgos en modelos de inteligencia artificial aplicados al reclutamiento mediante currículums sintéticos}

\begin{document}
\begin{titlepage}
\centering
\vspace*{1.2cm}
{\Large\bfseries Máster en Inteligencia Artificial\par}
\vspace{2.1cm}
{\large Trabajo Fin de Máster\par}
\vspace{1cm}
{\LARGE\bfseries\tituloTFM\par}
\vfill
{\large\textbf{Alumno}\par}
\vspace{0.2cm}
{\Large Jordan Rajo Vázquez\par}
\vspace{1.2cm}
{\large\textbf{Fecha de entrega}\par}
\vspace{0.2cm}
{\large 30 de octubre de 2026\par}
\vspace*{1cm}
\end{titlepage}

\pagenumbering{roman}
\chapter*{Resumen}
\addcontentsline{toc}{chapter}{Resumen}
Este trabajo estudia cómo el desequilibrio de representación y las etiquetas históricas sesgadas afectan a clasificadores de currículums en español. Se construye un banco de pruebas de cinco corpus independientes, 10.000 perfiles profesionales y 60.000 CVs principales, con seis variantes de género y edad por perfil. Una rúbrica explícita permite distinguir idoneidad sintética y decisión observada. Se comparan regresión logística y SVM lineal sobre TF--IDF, cinco escenarios de datos y seis tratamientos, incluidos limpieza de señales indirectas, reponderación y reparación simulada de una muestra de etiquetas. La evaluación separa rendimiento, errores interseccionales e invariancia emparejada e incorpora cambio de formato y sensibilidad a la rúbrica. La contribución es un protocolo inspeccionable con predicciones individuales, verificación independiente y material reproducible. La limpieza exhaustiva garantiza invariancia en este generador, pero no corrige por sí sola todas las decisiones erróneas ni acredita equidad sobre candidaturas reales. Los resultados se interpretan como evidencia sobre mecanismos controlados, no como estimaciones de discriminación del mercado laboral.

\noindent\textbf{Palabras clave:} selección de personal; equidad algorítmica; CVs sintéticos; procesamiento del lenguaje natural; auditoría; reproducibilidad.

\section*{Abstract}
This dissertation investigates how representation imbalance and biased training labels affect Spanish résumé classifiers. Five independently generated corpora contain 10,000 professional profiles and 60,000 primary résumés, with six paired gender--age variants per profile. An explicit rubric separates synthetic qualification from observed decisions. Logistic regression and linear SVM models using TF--IDF are evaluated under five data scenarios and six treatments, including proxy removal, reweighting and simulated partial label repair. Evaluation covers predictive performance, intersectional error rates, paired invariance, format shifts and rubric sensitivity. The contribution is an inspectable protocol with individual predictions, independent verification and reproducible materials. Exhaustive cleaning ensures invariance within the generator but does not guarantee correct classification or fairness for real applicants. Findings concern controlled mechanisms rather than the prevalence of labour-market discrimination.

\noindent\textbf{Keywords:} recruitment; algorithmic fairness; synthetic résumés; natural language processing; auditing; reproducibility.
\clearpage
\tableofcontents
\clearpage\listoffigures
\clearpage
\pagenumbering{arabic}

\chapter{Introducción}
\label{cap:introduccion}

\section{Contexto: inteligencia artificial y selección de personal}

La selección de personal combina la identificación de necesidades organizativas, la evaluación de candidaturas y la toma de decisiones sobre el acceso al empleo. Dentro de este proceso, el análisis de currículums constituye una primera aproximación a los conocimientos y la trayectoria profesional de una persona. La incorporación de herramientas de inteligencia artificial (IA) permite automatizar parte de ese análisis mediante la extracción de información, la clasificación de solicitudes o la ordenación de perfiles según su correspondencia con una oferta. Sin embargo, la transformación de un documento en una puntuación exige decidir qué información resulta relevante y qué criterio define una candidatura adecuada.

Estas decisiones no son exclusivamente técnicas. Un modelo entrenado para reproducir contrataciones anteriores aprende a aproximar el comportamiento registrado en esos datos, que no tiene por qué coincidir con una evaluación justa de las competencias. Raghavan et al.\ \cite{raghavan2020} examinan las prácticas de proveedores de herramientas de contratación y muestran la importancia de analizar conjuntamente la definición del objetivo, la validación y las afirmaciones sobre mitigación de sesgos. Su estudio sitúa el problema en el diseño completo del proceso, más allá de la elección de un algoritmo concreto.

El procesamiento del lenguaje natural (PLN) añade otra dimensión a esta cuestión. Un CV contiene cualificaciones explícitas, pero también nombres, fechas, formas de expresión y referencias contextuales. Parte de esa información puede actuar como señal de características personales ajenas a los requisitos del puesto. Por ello, un sistema puede producir diferencias entre grupos incluso cuando no recibe una variable estructurada denominada género o edad. La cuestión central de este TFM consiste en estudiar cómo se manifiestan esas diferencias en un entorno controlado y qué intervenciones permiten reducirlas.

\section{Justificación del problema y relevancia del estudio}

El acceso a una entrevista representa una oportunidad concreta. Un falso negativo puede excluir a una persona que satisface los requisitos de una oferta; si ese error se concentra en un grupo, la calidad global del modelo no describe adecuadamente sus consecuencias. Por ejemplo, dos clasificadores con el mismo porcentaje total de aciertos pueden distribuir los rechazos erróneos de forma muy distinta. Una evaluación útil debe considerar tanto la capacidad de identificar perfiles adecuados como el reparto de los errores.

La utilización de CVs sintéticos responde a la necesidad de construir comparaciones reproducibles. Permite fijar competencias y experiencia, variar atributos personales y conocer las reglas que originan las etiquetas. También facilita diseñar escenarios de entrenamiento con distintos grados de desequilibrio o de asociación entre atributos sensibles y decisiones favorables. El interés científico reside en observar la respuesta de los modelos a esas condiciones, distinguiendo los efectos introducidos deliberadamente de las diferencias que aparecen por la representación del texto o el aprendizaje.

Esta elección no convierte los datos sintéticos en una reproducción fiel del mercado laboral. Sus propiedades dependen del generador, las plantillas y las reglas utilizadas. Por tanto, el estudio se plantea como un banco de pruebas experimental: permite analizar mecanismos y comparar intervenciones dentro de los escenarios definidos, sin estimar directamente la prevalencia de discriminación en empresas reales. Esta delimitación hace viable una investigación de alcance acotado y evita atribuir a los resultados una generalidad que el diseño no puede sostener.

\section{Objetivo general y objetivos específicos}

El objetivo general es evaluar cómo modelos de IA entrenados con CVs sintéticos pueden reproducir o amplificar diferencias asociadas al género y la edad en la preselección de candidaturas, y comparar mecanismos de mitigación considerando simultáneamente la equidad y el rendimiento predictivo.

Este objetivo se concreta en las siguientes actuaciones:
\begin{enumerate}
\item Diseñar un conjunto de CVs sintéticos con variables profesionales, atributos personales y reglas de adecuación a una oferta explícitamente documentadas.
\item Construir escenarios de entrenamiento que permitan diferenciar el desequilibrio de representación del sesgo introducido en las etiquetas.
\item Entrenar y comparar una regresión logística y una SVM lineal con una representación TF--IDF común de los CVs.
\item Medir el rendimiento global y por grupos, así como las diferencias de selección y de errores por género, edad y sus intersecciones.
\item Analizar la sensibilidad de las predicciones mediante pares de CVs equivalentes en sus requisitos profesionales.
\item Evaluar cinco intervenciones de mitigación, además de la referencia original, y formular recomendaciones ajustadas a sus beneficios, costes y limitaciones.
\end{enumerate}

\section{Preguntas de investigación}

La investigación se organiza alrededor de cuatro preguntas. \textbf{PI1:} ¿en qué medida las diferencias de representación y de etiquetado del entrenamiento se trasladan a las predicciones? \textbf{PI2:} ¿cómo varían las disparidades por género y edad entre una regresión logística y una SVM lineal con la misma representación léxica del CV? \textbf{PI3:} ¿cambia la decisión cuando se modifican señales personales manteniendo constantes los requisitos profesionales? \textbf{PI4:} ¿qué mejoras de equidad se obtienen mediante las técnicas de mitigación y qué efecto tienen sobre la identificación de candidaturas adecuadas?

Estas preguntas no presuponen que todos los modelos discriminen ni que una técnica sea superior en cualquier escenario. La existencia, magnitud y dirección de las diferencias se evaluaron mediante los experimentos. También se considera un resultado relevante que una disparidad desaparezca bajo determinadas condiciones o que una intervención mejore una métrica y empeore otra.

\section{Alcance y enfoque del trabajo}

La tarea realizada fue la clasificación de candidaturas para una oferta ficticia de analista de datos. La salida favorable representa pasar a una fase posterior de evaluación, no recibir una oferta de empleo. Se utilizaron CVs textuales en español y se estudiaron género y edad. Otros atributos quedan fuera del experimento y pueden considerarse en ampliaciones posteriores.

El diseño comparó regresión logística y SVM lineal sobre TF--IDF, con entrenamiento supervisado a partir de datos sintéticos. Se compararon reponderación, supresión de cabeceras, limpieza de señales indirectas y combinaciones con reparación parcial de etiquetas. Se realizaron pruebas de sensibilidad y evaluación con formatos reservados. El capítulo metodológico detalla modelos, tamaños, particiones e hiperparámetros. Las representaciones semánticas preentrenadas se mantienen como antecedente teórico y posible ampliación; no se utilizaron en esta ejecución.

Las etiquetas de evaluación se definieron mediante una regla profesional de referencia independiente de los atributos sensibles. Al introducir sesgo en el entrenamiento, se conservó esa referencia para distinguir la imitación de decisiones sesgadas de la adecuación profesional simulada. La regla sigue siendo una construcción del estudio y no una medida demostrada del desempeño laboral futuro.



\chapter{Marco teórico y estado del arte}
\label{cap:marco}

\section{Aprendizaje automático aplicado al reclutamiento}

\subsection{De los requisitos profesionales a una tarea predictiva}

El filtrado de CVs puede formularse como clasificación, estimación de una puntuación u ordenación de candidatos. Estas formulaciones producen decisiones diferentes: un clasificador compara cada perfil con un umbral, mientras que un sistema de ordenación distribuye posiciones dentro de una lista. En este TFM se adopta la clasificación binaria para facilitar el análisis de errores y comparaciones entre grupos. Si se añadiera una selección de los primeros puestos de una lista, sería necesario evaluar además cómo influye el tamaño de la lista en la exposición de cada grupo.

Sea $X$ la representación de las características profesionales y del texto de un CV, $A$ el atributo personal estudiado, $Y^{*}$ la adecuación definida por la regla de referencia y $\widetilde{Y}$ la etiqueta utilizada para entrenar. Un modelo produce una puntuación $s_{\theta}(X)$ y una decisión:
\begin{equation}
\widehat{Y}=\mathbb{I}\{s_{\theta}(X)\geq t\},
\end{equation}
donde $\theta$ recoge los parámetros aprendidos, $t$ es el umbral y $\mathbb{I}$ la función indicadora. Aunque $A$ no figure como entrada explícita, su información puede estar contenida indirectamente en $X$. En los escenarios sin alteración de etiquetas se tendrá $\widetilde{Y}=Y^{*}$; en los escenarios sesgados, ambas variables podrán diferir.

Esta distinción evita una ambigüedad central del experimento. Un modelo puede aprender correctamente las etiquetas de entrenamiento y, al mismo tiempo, alejarse de la regla profesional prevista. Por ello, reproducir el etiquetado observado no bastará para concluir que el sistema es adecuado. La validación deberá especificar siempre qué variable se utiliza como referencia y qué significado tiene una decisión favorable.

\subsection{Representaciones textuales y modelos de referencia}

Una representación TF--IDF asigna peso a los términos según su frecuencia en un documento y su frecuencia inversa en la colección. Así, combina información local con la capacidad discriminativa de una palabra dentro del corpus \cite{manning2008}. Asociada a una regresión logística, ofrece un punto de partida de coste computacional reducido y permite inspeccionar la contribución de características léxicas. Su limitación principal para este caso es que expresiones distintas de una misma competencia pueden tratarse como características separadas.

Las representaciones semánticas densas buscan aproximar semejanzas de significado. Sentence-BERT propone obtener vectores de frases mediante una arquitectura entrenada para comparaciones semánticas \cite{reimers2019}. Una ampliación con un codificador apropiado para español deberá comprobar la longitud admitida, la estrategia de agregación de secciones y la cobertura lingüística. La idoneidad de un modelo para similitud textual no demuestra, por sí misma, su validez para juzgar adecuación profesional.

En el experimento realizado se mantuvo TF--IDF como representación común y se compararon dos clasificadores. Esta elección permite estudiar el efecto de la familia de aprendizaje bajo las mismas particiones, etiquetas de referencia y búsqueda de regularización. La comparación con representaciones semánticas queda fuera del alcance empírico.

\section{Origen y manifestaciones del sesgo}

\subsection{Sesgos a lo largo del ciclo de vida}

Suresh y Guttag \cite{suresh2021} distinguen fuentes de daño relacionadas con la historia de los datos, su representación y medición, la agregación, el aprendizaje, la evaluación y el despliegue. Su marco permite localizar problemas que quedarían ocultos si el sesgo se atribuyera únicamente al conjunto de entrenamiento. Aplicado al reclutamiento, invita a examinar quién aparece en los datos, qué se considera éxito y cómo se utiliza la predicción.

En este estudio, el \textit{sesgo de representación} se operacionalizará mediante diferencias en el número de CVs disponibles por grupo. El \textit{sesgo de etiquetado} se simulará alterando la probabilidad de recibir una etiqueta favorable para perfiles de cualificación comparable. Se tratarán como factores distintos: aumentar la presencia de un grupo no garantiza que sus etiquetas sean adecuadas. Asimismo, el generador se revisará para evitar que asigne sistemáticamente un estilo, una longitud o determinadas competencias a una categoría demográfica.

La evaluación agregada puede ocultar resultados adversos en subgrupos pequeños. En consecuencia, cada métrica se acompañará del número de observaciones que la sustenta. Si un subgrupo carece de casos positivos o negativos, algunas tasas no estarán definidas; no se sustituirán por cero, porque ello transmitiría una falsa ausencia de errores.

\subsection{Atributos explícitos, variables indirectas e intersecciones}

La eliminación de nombres y pronombres no asegura que desaparezca la información asociada al género. De-Arteaga et al.\ \cite{dearteaga2019}, en una tarea de clasificación de ocupaciones a partir de biografías, observaron diferencias de rendimiento y señales indirectas incluso después de retirar indicadores explícitos. Este antecedente justifica comprobar empíricamente la eficacia de la supresión de atributos. Su tarea, sin embargo, es clasificar ocupaciones y no decidir contrataciones, por lo que sus resultados no deben trasladarse directamente a la selección de personal.

En el experimento propuesto, los nombres se tratarán como señales lingüísticas manipuladas y no como un procedimiento para inferir la identidad real de una persona. Las categorías de género utilizadas se documentan como decisiones de diseño. Si el tamaño muestral obliga a limitar las categorías comparadas, esa restricción se explicitará y no se presentará como una descripción exhaustiva del género.

La edad plantea una dificultad adicional: suele estar relacionada con la duración de la experiencia y con fechas educativas. Cambiar únicamente una fecha de nacimiento puede producir un CV incoherente. Para evitarlo se diseñarán perfiles plausibles con experiencia relevante comparable y se distinguirá la modificación de una señal explícita del estudio de trayectorias completas. El análisis interseccional cruzará género y grupos de edad para comprobar si los promedios marginales ocultan patrones específicos. El experimento utiliza las edades puntuales de 28, 43 y 58 años, fijadas antes de observar resultados; no estima efectos para intervalos completos.

\section{Evaluación del rendimiento y de la equidad}

\subsection{Calidad predictiva y distribución de errores}

La matriz de confusión distinguirá verdaderos positivos, falsos positivos, verdaderos negativos y falsos negativos respecto de $Y^{*}$. En esta aplicación, un falso negativo será un CV que cumple la regla de referencia y queda excluido. Se calcularán precisión, sensibilidad y F1, junto con resultados desagregados. La exactitud global se interpretará atendiendo a la distribución de clases: un modelo que rechaza casi todas las candidaturas puede acertar con frecuencia cuando la clase favorable es minoritaria y resultar poco útil para encontrar perfiles adecuados.

El umbral forma parte del sistema evaluado. Cambiarlo modifica tanto los errores como las tasas de selección, por lo que deberá quedar especificado antes de evaluar la prueba. En este experimento se utilizaron los umbrales convencionales fijos de cada clasificador. El conjunto de prueba se reservará para la comparación final. Las variantes de un mismo perfil permanecerán en una única partición para evitar que el modelo encuentre en la evaluación textos casi idénticos a los utilizados durante el entrenamiento.

\subsection{Paridad demográfica e igualdad de oportunidades}

Para un grupo $a$, se define la tasa de selección $SR_a=P(\widehat{Y}=1\mid A=a)$. Una medida descriptiva de disparidad es:
\begin{equation}
\Delta_{\mathrm{DP}}=\max_a SR_a-\min_a SR_a.
\end{equation}
Un valor nulo representa paridad de selección entre los grupos considerados. Esta medida no utiliza la cualificación de referencia, de modo que necesita interpretarse junto con la composición del conjunto evaluado. En el banco de pruebas se podrán construir grupos con cualificación equivalente para facilitar esa interpretación; no se supondrá que tal equivalencia existe automáticamente en cualquier población.

Hardt et al.\ \cite{hardt2016} formalizan la igualdad de oportunidades mediante la igualdad de tasas de verdaderos positivos y la igualdad de probabilidades condicionadas mediante la coincidencia de verdaderos y falsos positivos. Con $TPR_a=P(\widehat{Y}=1\mid Y^{*}=1,A=a)$ y $FPR_a=P(\widehat{Y}=1\mid Y^{*}=0,A=a)$, se utilizarán:
\begin{align}
\Delta_{\mathrm{EOpp}}&=\max_a TPR_a-\min_a TPR_a,\\
\Delta_{\mathrm{EOdds}}&=\max\left\{\Delta_{\mathrm{EOpp}},\max_a FPR_a-\min_a FPR_a\right\}.
\end{align}
La primera diferencia permite estudiar si los CVs adecuados tienen oportunidades comparables de superar el filtro. La segunda incorpora también la selección incorrecta. Ambas dependen de la validez de $Y^{*}$: si la referencia incorpora un criterio injustificado, una aparente igualdad de errores puede ser engañosa.

\subsection{Compatibilidad de criterios e incertidumbre}

Las nociones de equidad responden a preguntas diferentes. Kleinberg et al.\ \cite{kleinberg2017} demuestran incompatibilidades entre condiciones de calibración y equilibrio de puntuaciones entre grupos, salvo casos restringidos. Su resultado no implica que toda mejora de equidad exija perder exactitud, sino que no cabe asumir que distintos criterios puedan satisfacerse simultáneamente en cualquier distribución.

En este TFM se priorizará la igualdad de oportunidades respecto de la regla profesional, y se utilizarán las restantes medidas como información complementaria. La comparación recogerá cambios absolutos, tamaños de grupo y variabilidad entre repeticiones. Se plantearán intervalos mediante remuestreo de perfiles base, conservando juntas sus variantes. Las diferencias pequeñas no se interpretan automáticamente como ausencia de sesgo, ni la significación estadística sustituirá la valoración de su magnitud práctica.

\section{Pruebas emparejadas y datos sintéticos}

\subsection{Comparaciones controladas de CVs}

Kusner et al.\ \cite{kusner2017} desarrollan una noción causal de equidad contrafactual que compara decisiones en mundos alternativos definidos mediante un modelo causal. Esa noción es más exigente que sustituir un nombre en un documento. En este trabajo, las pruebas se denominarán \textit{comparaciones emparejadas de sensibilidad}: observarán qué sucede al cambiar señales previamente especificadas, sin afirmar que identifican todos los efectos causales del género o la edad.

Para $n$ pares de CVs equivalentes en los requisitos relevantes, una medida utilizada es la diferencia media absoluta de puntuación:
\begin{equation}
D_{\mathrm{par}}=\frac{1}{n}\sum_{i=1}^{n}\left|s_{\theta}(x_i^{(a)})-s_{\theta}(x_i^{(b)})\right|.
\end{equation}
Se informará también de la diferencia con signo y de la proporción de pares cuya decisión cambia. La primera evita que efectos opuestos se cancelen; la segunda indica qué variante resulta favorecida; la tercera conecta el cambio de puntuación con la preselección efectiva. Estas medidas describirán exclusivamente las modificaciones ensayadas.

\subsection{Ventajas y límites de la generación sintética}

Saldivar et al.\ \cite{saldivar2025} presentan una propuesta de construcción de CVs sintéticos basada en características de materiales obtenidos mediante donación de datos, junto con una colección de 1.730 CVs. Este preprint constituye un antecedente directamente relacionado con la disponibilidad de recursos para evaluar herramientas de contratación. Su aportación apoya la utilidad de documentar el proceso generativo, pero no demuestra que cualquier conjunto sintético reproduzca la diversidad del empleo real.

Para este TFM, la generación se organizará en dos niveles: una ficha estructurada que establezca las competencias y una realización textual que las exprese. La etiqueta de referencia se calculará a partir de la ficha, evitando que dependa de la fluidez de la redacción. Las plantillas se distribuirán entre grupos y se revisarán duplicados, contradicciones y asociaciones accidentales. En la ejecución realizada se utilizaron reglas y plantillas, sin apoyo de un modelo generativo.

La separación entre generación, etiquetado y evaluación permitirá distinguir tareas que, de otro modo, podrían confundirse. Un sistema que identifica la plantilla asociada a la clase positiva no ha demostrado comprender requisitos profesionales. Asimismo, evaluar solo con el mismo estilo de texto utilizado para entrenar puede favorecer una generalización aparente. La evaluación con estilos de redacción reservados constituiría una comprobación adicional. No se realizó en este experimento, cuyas particiones comparten el mismo repertorio de plantillas.

\section{Técnicas de mitigación de sesgos}

\subsection{Intervenciones sobre los datos}

Kamiran y Calders \cite{kamiran2012} estudian procedimientos de preprocesamiento como supresión de atributos, modificación de etiquetas, reponderación y remuestreo. La reponderación modifica la contribución de los ejemplos durante el aprendizaje para reducir asociaciones no deseadas entre grupo y etiqueta. Su atractivo experimental es que permite mantener la familia de modelos y comparar el efecto de cambiar los pesos de entrenamiento.

La primera intervención repondera combinaciones de grupo y etiqueta únicamente dentro del entrenamiento. Al interpretarla deben considerarse pesos extremos y combinaciones escasas, porque un grupo con muy pocos ejemplos no adquiere diversidad por recibir más peso. La segunda intervención elimina la cabecera demográfica mediante reglas documentadas. La comparación antes y después de la transformación permite observar cómo cambian el rendimiento y las diferencias entre grupos.

Estas intervenciones abordarán mecanismos distintos. La reponderación actúa sobre la influencia relativa de los ejemplos; la supresión modifica la información disponible. Ninguna se considera satisfactoria por su intención: ambas deberán mostrar mejoras sobre el mismo conjunto de prueba y en las métricas predefinidas. Además, los atributos de auditoría se conservarán separados de las entradas del clasificador para poder medir resultados por grupo.

\subsection{Restricciones durante el aprendizaje y ajustes posteriores}

Agarwal et al.\ \cite{agarwal2018} plantean la clasificación con restricciones de equidad como una secuencia de problemas de aprendizaje sensibles al coste. Este enfoque permite incorporar explícitamente condiciones estadísticas durante el entrenamiento. Constituye una alternativa a modificar los datos, aunque requiere elegir la restricción y evaluar su cumplimiento fuera de la muestra de ajuste.

El posprocesamiento ajusta decisiones de un predictor ya entrenado. Puede modificar umbrales para aproximar un criterio de equidad, como muestran Hardt et al.\ \cite{hardt2016}. Sin embargo, algunas variantes requieren conocer el grupo en el momento de decidir. Por ello, el análisis de una técnica debe considerar no solo sus métricas, sino también qué información necesita para aplicarse. En el alcance inicial del TFM, estas familias servirán de contexto y posible ampliación; la comparación principal se concentrará en las dos intervenciones sobre datos descritas anteriormente.

\section{Consideraciones éticas y marco regulatorio}

\subsection{Justificación, transparencia y supervisión}

La evaluación ética requiere justificar por qué una característica interviene en la decisión y cómo se relaciona con el puesto. En el experimento, esta exigencia se traduce en una regla explícita de adecuación, una descripción de los errores relevantes y una separación entre observación y recomendación. Una explicación legible del modelo puede ayudar a inspeccionar una predicción, pero no convierte por sí sola el criterio aprendido en aceptable.

Las propuestas de \textit{datasheets} para conjuntos de datos \cite{gebru2021} y de \textit{model cards} para modelos \cite{mitchell2019} ofrecen referencias para documentar procedencia, composición, usos previstos y resultados desagregados. El TFM adapta esos principios a una ficha del conjunto sintético y a un registro de cada modelo. También documenta las limitaciones de las categorías demográficas y las condiciones en las que no puede interpretarse una métrica.

La supervisión humana se considera parte del proceso de uso. Una puntuación no debería sustituir la revisión de los requisitos y de la información del CV. Como el trabajo no despliega un sistema con personas candidatas reales, las recomendaciones de supervisión tienen carácter de diseño y no se presentan como medidas cuya eficacia se haya comprobado experimentalmente.

\subsection{Normativa europea y española}

El Reglamento (UE) 2024/1689 incluye en su anexo III determinados sistemas destinados a la contratación, al filtrado de solicitudes y a la evaluación de candidatos. La clasificación jurídica concreta requiere atender a su finalidad y a las condiciones del artículo 6. El marco de alto riesgo contempla aspectos como gestión de riesgos, gobernanza de datos, documentación y supervisión humana \cite{aiact2024}. La información oficial de la Comisión consultada el 13 de octubre de 2026 sitúa la aplicación de las reglas de los ámbitos de alto riesgo del anexo III, incluido empleo, el 2 de diciembre de 2027, tras la modificación del calendario \cite{comision2026,omnibus2026}. Se distingue así la existencia del marco de la fecha de aplicación de sus obligaciones específicas.

El Reglamento General de Protección de Datos resulta relevante cuando se tratan datos personales. Su artículo 22 regula las decisiones basadas únicamente en tratamiento automatizado que produzcan efectos jurídicos o afecten significativamente de modo similar, con excepciones y salvaguardias. No equivale a una prohibición general de utilizar IA en selección. Además, el uso de información debe respetar principios como finalidad, minimización y licitud \cite{rgpd2016}. Los atributos sensibles para una auditoría de equidad no coinciden necesariamente con las categorías especiales del artículo 9; género y edad no deben agruparse automáticamente bajo esa categoría jurídica.

En España, la Ley 15/2022 aborda la igualdad de trato y la no discriminación, incluido el empleo. Su artículo 23 recoge previsiones sobre algoritmos e IA, con referencias a sesgos, transparencia y rendición de cuentas \cite{ley152022}. En este TFM, estas normas sustentan la relevancia del problema y la necesidad de documentación. El cumplimiento de una métrica estadística no se interpreta como una certificación de cumplimiento legal, y el experimento académico no se equipara a la puesta en servicio de una herramienta comercial.

\section{Estado del arte y síntesis crítica}
La revisión se organizó alrededor de tres preguntas: cómo se genera el sesgo en selección automatizada, cómo se auditan diferencias entre candidaturas comparables y qué información necesita una mitigación. Se priorizaron publicaciones originales, documentación técnica oficial y normas primarias. Se combinaron términos relativos a \emph{algorithmic hiring}, \emph{resume screening}, \emph{synthetic CVs}, \emph{fairness} y \emph{audit studies}. Se incluyeron estudios con tarea, datos o mecanismo identificables y se distinguió entre clasificación de profesiones, recuperación de CVs y decisiones de selección. No se presenta un recuento de cribado ni una evaluación exhaustiva de calidad, porque la revisión es narrativa y no sigue un protocolo sistemático registrado.


\subsection{De las representaciones profesionales a las auditorías de selección}

Los antecedentes seleccionados recorren tres niveles complementarios. Raghavan et al.\ \cite{raghavan2020} examinan las afirmaciones y prácticas organizativas de las herramientas de contratación; De-Arteaga et al.\ \cite{dearteaga2019} estudian cómo las representaciones de biografías mantienen diferencias asociadas al género; y Wilson y Caliskan \cite{wilson2024} analizan una simulación de recuperación de CVs mediante modelos de representaciones textuales. Este último trabajo muestra la pertinencia de auditar señales asociadas a nombres e intersecciones demográficas y de considerar factores como la longitud del documento. Aquí se toma como antecedente de diseño, sin trasladar sus porcentajes a otro idioma, población o modelo.

La evidencia sobre datos sintéticos complementa esas auditorías. La propuesta de Saldivar et al.\ \cite{saldivar2025} aporta un recurso orientado a estudiar discriminación en contratación, aunque su condición de preprint y las propiedades de su generador deben tenerse presentes. El TFM se apoya en esta línea para construir un experimento propio, sin asumir que la generación sintética resuelve automáticamente los problemas de validez.

Una aportación reciente especialmente relevante para la edad es la de Sariola et al.\ \cite{sariola2026}. A partir de datos de una auditoría de discriminación por edad, evalúan intervenciones de remuestreo y muestran que una aparente paridad bajo medidas convencionales puede coexistir con disparidades al utilizar una evaluación más adecuada. La implicación para este trabajo es concreta: la mitigación no debe juzgarse únicamente con las mismas etiquetas y supuestos que pudieron originar el sesgo.

\subsection{Comparación de aportaciones y límites}

La tabla~\ref{tab:antecedentes} sintetiza las diferencias entre los antecedentes directamente relacionados con el caso. La selección combina trabajos fundacionales y contribuciones recientes localizadas en fuentes académicas primarias. Se trata de una revisión narrativa focalizada, no de una revisión sistemática con pretensión de exhaustividad. Las fuentes y el marco normativo se han consultado hasta el 5 de octubre de 2026, con actualización del calendario regulatorio el 13 de octubre de 2026.

\begin{table}[htbp]
\centering
\small
\setlength{\tabcolsep}{4pt}
\renewcommand{\arraystretch}{1.12}
\caption{Antecedentes relevantes y relación con el diseño del TFM. Síntesis propia.}
\label{tab:antecedentes}
\begin{tabularx}{\textwidth}{@{}p{2.7cm}XX@{}}
\toprule
\textbf{Trabajo} & \textbf{Aportación al estudio} & \textbf{Límite de transferencia}\\
\midrule
Raghavan et al.\ (2020) & Conecta validación, objetivo y prácticas de mitigación. & No estima el sesgo de los modelos entrenados aquí.\\
De-Arteaga et al.\ (2019) & Motiva examinar señales indirectas en texto. & Clasifica ocupaciones; no evalúa contrataciones.\\
Wilson y Caliskan (2024) & Orienta auditorías de CVs y análisis interseccionales. & Estudia recuperación textual en un contexto específico.\\
Saldivar et al.\ (2025) & Aporta un precedente de CVs sintéticos documentados. & La validez depende del generador y la población representada.\\
Sariola et al.\ (2026) & Cuestiona evaluar mitigaciones con referencias sesgadas. & Utiliza un diseño y datos de auditoría distintos.\\
\bottomrule
\end{tabularx}
\end{table}

\subsection{Posicionamiento y contribución del TFM}

El espacio de contribución del trabajo consiste en integrar un diseño manejable y reproducible: CVs sintéticos en español, dos clasificadores con representación textual común, género y edad como ejes principales, y comparación antes y después de mitigar. Su interés no depende de proponer un algoritmo nuevo, sino de hacer explícitas las relaciones entre generación de datos, etiquetas, decisiones y métricas.

La aportación es un protocolo que permite identificar bajo qué escenarios aparece una disparidad y qué cambia al intervenir. La referencia profesional independiente, la reserva de perfiles completos por partición y las pruebas emparejadas refuerzan la interpretación de los resultados. Aun así, permanecen límites derivados de la simplificación de trayectorias, el tamaño de los subgrupos y la ausencia de resultados laborales reales. Las conclusiones se formulan dentro de ese alcance y señalan qué comprobaciones adicionales serían necesarias antes de trasladar una recomendación a una organización.

